\documentclass[a4paper, serif, 10pt]{moderncv}

%% moderncv
% style and color
\moderncvstyle{banking}
\moderncvcolor{black}

% renew moderncv command to adapt for CJK colon
\renewcommand*{\cvitem}[3][.25em]{%
  \ifstrempty{#2}{}{\hintstyle{#2}：}{#3}%
  \par\addvspace{#1}}

\renewcommand*{\cvitemwithcomment}[4][.25em]{%
  \savebox{\cvitemwithcommentmainbox}{\ifstrempty{#2}{}{\hintstyle{#2}：}#3}%
  \setlength{\cvitemwithcommentmainlength}{\widthof{\usebox{\cvitemwithcommentmainbox}}}%
  \setlength{\cvitemwithcommentcommentlength}{\maincolumnwidth-\separatorcolumnwidth-\cvitemwithcommentmainlength}%
  \begin{minipage}[t]{\cvitemwithcommentmainlength}\usebox{\cvitemwithcommentmainbox}\end{minipage}%
  \hfill% fill of \separatorcolumnwidth
  \begin{minipage}[t]{\cvitemwithcommentcommentlength}\raggedleft\small\itshape#4\end{minipage}%
  \par\addvspace{#1}}

%% page layout/margins
\usepackage[top=2.5cm, bottom=2.5cm, left=1.5cm, right=1.5cm]{geometry}

\nopagenumbers{}

%% Babel config for English language
\usepackage[english]{babel}

%% fontspec
\usepackage{fontspec}

\IfFontExistsTF{Linux Libertine}{
  \setmainfont[Ligatures={TeX, Common}, Numbers=Lining, ItalicFont=Linux Libertine]{Linux Libertine}
}{}
\IfFontExistsTF{Linux Libertine O}{
  \setmainfont[Ligatures={TeX, Common}, Numbers=Lining, ItalicFont=Linux Libertine O]{Linux Libertine O}
}{}

%% CTeX
% CJK support, used to show CJK characters in the resume
%
% - fontset=none: disable builtin fontset but instead set the CJK font manually
% - heading=false: disable ctex heading
% - punct=kaiming: use kaiming punctuations styles for CJK
% - scheme=plain: use plain scheme, do not override \normalsize font size
% - space=auto: space settings for CJK characters
%
% ref:
% - http://ctan.mirrorcatalogs.com/language/chinese/ctex/ctex.pdf
\usepackage[UTF8, heading=false, punct=kaiming, scheme=plain, space=auto]{ctex}

\IfFontExistsTF{Noto Serif CJK SC}{
  \setCJKmainfont{Noto Serif CJK SC}
}{}
\IfFontExistsTF{Noto Sans CJK SC}{
  \setCJKsansfont{Noto Sans CJK SC}
}{}

%% line spacing
\usepackage{setspace}
\setstretch{1.125}

%% URL styling - use normal text instead of monospace
\urlstyle{same}

% Auto-underline all links
% Defer until \begin{document} so hyperref is already loaded
\AtBeginDocument{%
  \let\oldhref\href
  \renewcommand{\href}[2]{\oldhref{#1}{\underline{#2}}}%
}

%% Patch \httplink and \httpslink to handle full URLs
% The original moderncv macros always prepend http:// or https:// to URLs.
% This causes issues when the URL already has a protocol (e.g., https://example.com)
% resulting in malformed URLs like https://https://example.com.
% This patch checks if the URL already contains "://" and uses it directly if so.
% Note: We use \str_set:Nx (x-expansion) to expand macros like \@homepage first.
\ExplSyntaxOn
\str_new:N \g_yamlresume_url_str

\RenewDocumentCommand{\httpslink}{O{}m}{
  \str_set:Nx \g_yamlresume_url_str {#2}
  \str_if_in:NnTF \g_yamlresume_url_str {://}
    { \url{#2} }
    { \url{https://#2} }
}

\RenewDocumentCommand{\httplink}{O{}m}{
  \str_set:Nx \g_yamlresume_url_str {#2}
  \str_if_in:NnTF \g_yamlresume_url_str {://}
    { \url{#2} }
    { \url{http://#2} }
}
\ExplSyntaxOff

\name{林雅婷}{}
\title{資深客戶成功經理 | 專注於 SaaS 與企業軟體}
\phone[mobile]{+886 912 345 678}
\email{yating.lin@gmail.com}
\homepage{https://www.linkedin.com/in/yatinglin}

\address{中國臺灣，信義區，台北市，台北市信義區松高路 88 號 12 樓，110}{}{}

\extrainfo{{\small \faLinkedin}\ \href{https://www.linkedin.com/in/yatinglin}{@yatinglin} {} {} {} • {} {} {} 
{\small \faGithub}\ \href{https://github.com/yatinglin}{@yatinglin} {} {} {} • {} {} {} 
{\small \faMedium}\ \href{https://medium.com/@yatinglin}{@yatinglin}}

\begin{document}

\maketitle

\section{簡介}

\cvline{}{具備 6 年以上客戶成功與帳戶管理經驗，專精於 SaaS 與企業軟體產業。擅長透過數據驅動策略提升客戶續約率、擴展收入與滿意度，並帶領團隊達成組織目標。
%
\begin{itemize}
\item 管理超過 30 家中大型企業客戶，年度續約率維持 95\% 以上，淨收入留存率 (NRR) 達 118\%
\item 建立客戶健康度評分模型與主動關懷流程，成功降低客戶流失風險 40\%
\item 領導 5 人客戶成功團隊，制定 onboarding 標準作業程序，縮短新客戶上線時間 35\%
\item 熟悉 Salesforce、HubSpot、Tableau 等工具，具備 SQL 與數據分析能力
\item 卓越的跨部門溝通與談判技巧，能有效協調產品、工程與業務團隊
\end{itemize}}
\section{教育背景}

\cventry{2014年9月–2018年6月}
        {學士，企業管理學系，成績：3.9/4.3}
        {國立臺灣大學}
        {\url{https://www.ntu.edu.tw/}}
        {}
        {\begin{itemize}
\item 主修企業管理，輔修資訊管理，奠定商業策略與數據分析基礎
\item 擔任系學會公關長，籌辦多場企業參訪與講座，提升溝通與協調能力
\item 畢業專題「以客戶旅程地圖優化 SaaS onboarding 流程」獲教授推薦
\end{itemize}
\textbf{課程}：行銷管理、消費者行為、策略管理、人力資源管理、財務管理、商業分析、專案管理、談判理論}
\section{工作經歷}

\cventry{2022年1月至今}
        {資深客戶成功經理}
        {雲端脈動科技股份有限公司}
        {\url{https://www.cloudpulse.tw}}
        {}
        {\begin{itemize}
\item 負責管理 30+ 家中大型企業客戶，年度續約率維持 95\% 以上，淨收入留存率 (NRR) 達 118\%
\item 領導 5 人客戶成功團隊，制定客戶健康度評分模型與主動關懷流程，降低流失風險 40\%
\item 與產品團隊協作，根據客戶回饋推動 12 項功能優化，提升客戶滿意度 (CSAT) 至 4.8/5.0
\item 建立客戶成功手冊與 onboarding 標準作業程序，將新客戶上線時間縮短 35\%
\item 定期舉辦客戶成功工作坊與用戶社群活動，促進客戶間交流與產品採用
\end{itemize}
\textbf{關鍵字}：客戶關係管理、續約策略、數據分析、團隊領導、SaaS}

\cventry{2018年7月–2021年12月}
        {客戶成功經理}
        {智創軟體服務有限公司}
        {\url{https://www.smartsoft.com.tw}}
        {}
        {\begin{itemize}
\item 管理 50+ 家中小企業客戶，負責產品導入、教育訓練與日常諮詢，客戶續約率達 90\%
\item 導入 CRM 系統 (Salesforce) 並優化客戶歷程追蹤，提升團隊工作效率 25\%
\item 分析客戶使用數據，主動識別潛在流失客戶並制定挽回計畫，成功挽回 15\% 即將流失客戶
\item 擔任客戶與產品部門之間的橋樑，每季彙整客戶需求並優先級排序，協助產品路線圖規劃
\item 獲得 2020 年度最佳客戶成功經理獎項，表彰卓越客戶服務與跨部門協作
\end{itemize}
\textbf{關鍵字}：客戶導入、教育訓練、Salesforce、客戶留存、跨部門溝通}
\section{語言}

\cvline{中文}{母語或雙語水平 \hfill \textbf{關鍵字}：繁體中文、台語}
\cvline{英語}{完全專業水平 \hfill \textbf{關鍵字}：TOEIC 950、商業書信、國際客戶溝通}
\section{專業技能}

\cvline{客戶成功與關係管理}{專家 \hfill \textbf{關鍵字}：客戶生命週期管理、續約與擴展、客戶健康度分析、淨推薦值 (NPS)、客戶滿意度 (CSAT)、QBR 規劃}
\cvline{數據分析與工具}{高級 \hfill \textbf{關鍵字}：Salesforce、HubSpot、Tableau、Excel、SQL、Google Analytics}
\cvline{專案與團隊管理}{高級 \hfill \textbf{關鍵字}：敏捷開發、Jira、Confluence、跨部門協作、團隊領導、簡報技巧}
\cvline{語言與溝通}{專家 \hfill \textbf{關鍵字}：中文 (母語)、英文 (專業)、談判、客戶簡報、衝突解決}
\section{榮譽}

\cventry{2020年12月}
        {智創軟體服務有限公司}
        {年度最佳客戶成功經理}
        {}
        {}
        {表彰在客戶續約率、客戶滿意度及跨部門協作方面的卓越表現，全年客戶留存率達 92\%，為團隊最高。}

\cventry{2023年6月}
        {雲端脈動科技股份有限公司}
        {卓越服務獎}
        {}
        {}
        {因成功輔導三家大型企業客戶完成數位轉型，客戶 NPS 達 75 分，獲公司季度卓越服務獎。}
\section{證書}

\cventry{2021年3月}
        {Salesforce}
        {Salesforce 認證管理員}
        {\url{https://trailhead.salesforce.com/en/credentials/administrator}}
        {}
        {}

\cventry{2022年8月}
        {Google}
        {Google 數據分析專業認證}
        {\url{https://grow.google/certificates/data-analytics/}}
        {}
        {}

\cventry{2019年11月}
        {PMI}
        {專案管理專業人士 (PMP) 認證}
        {\url{https://www.pmi.org/certifications/project-management-pmp}}
        {}
        {}
\section{出版刊物}

\cventry{2021年9月}
        {以客戶成功為核心的 SaaS 訂閱經濟策略}
        {台灣管理學刊}
        {\url{https://www.tamjournal.org.tw/article/12345}}
        {}
        {\begin{itemize}
\item 探討 SaaS 企業如何透過客戶成功管理提升續約率與客戶終身價值
\item 提出客戶健康度評分模型與主動關懷機制，並以台灣企業個案驗證
\item 獲選為當期最佳實務案例論文
\end{itemize}}
\section{引薦}

\cventry{\emaillink[ming.wang@cloudpulse.tw]{ming.wang@cloudpulse.tw}}
        {雲端脈動科技股份有限公司 營運副總}
        {王大明}
        {+886 2 1234 5678}
        {}
        {林雅婷在客戶成功領域表現卓越，具備數據驅動思維與卓越的溝通能力，能有效帶領團隊達成目標，是難得的人才。}
\section{項目}

\cventry{2022年1月–2022年6月}
        {基於 Python 與 Tableau 開發的客戶健康度監控系統}
        {客戶健康度儀表板}
        {\url{https://github.com/yatinglin/customer-health-dashboard}}
        {}
        {\begin{itemize}
\item 整合 Salesforce、產品使用數據與客服工單，計算客戶健康度分數
\item 提供即時警示與建議行動，協助客戶成功團隊主動管理風險客戶
\item 導入後客戶流失率降低 25\%，團隊工作效率提升 30\%
\end{itemize}
\textbf{關鍵字}：Python、Tableau、Salesforce API、數據視覺化}

\cventry{2023年3月–2023年8月}
        {設計並導入自動化 onboarding 流程，縮短客戶上線時間}
        {SaaS 客戶 onboarding 自動化流程}
        {\url{https://www.cloudpulse.tw/onboarding}}
        {}
        {\begin{itemize}
\item 使用 HubSpot 與 Zapier 建立自動化郵件序列與任務分派
\item 將新客戶上線時間從 14 天縮短至 9 天，提升客戶早期採用率
\item 專案成果獲選為公司年度最佳內部流程改善案例
\end{itemize}
\textbf{關鍵字}：HubSpot、Zapier、流程自動化、客戶體驗}
\section{興趣愛好}

\cvline{志工服務}{偏鄉教育、動物保護、社區營造}
\cvline{戶外活動}{登山、潛水、馬拉松}
\cvline{閱讀與寫作}{商業策略、客戶體驗、心理學}
\section{志願服務}

\cventry{2019年9月–2023年6月}
        {志工講師}
        {台灣偏鄉教育發展協會}
        {\url{https://www.ruraledu.tw}}
        {}
        {\begin{itemize}
\item 每月前往偏鄉國小教授商業基礎與職涯探索課程，累計服務 200+ 小時
\item 設計互動式教案，提升學生學習動機與自信心
\item 協助協會舉辦年度營隊，招募並培訓 30 位大專院校志工
\end{itemize}}
    
\end{document}